% -*- coding: utf-8 -*-

\chapter{总结与展望} \label{chpt:conclusion}

\section{工作总结}

% 本文针对深度神经网络面临的对抗样本威胁，
% 研究了基于显著性图的遮蔽攻击方法及其对抗性训练策略。
% 主要完成了以下工作：
%
% 1. 提出了基于梯度显著性图的遮蔽攻击方法
%    利用梯度显著性图定位对分类决策最关键的像素区域，
%    设计了固定遮蔽攻击和自适应遮蔽攻击两种方式。
%    与基于梯度积分的方法相比，本文方法仅需单次反向传播
%    即可完成关键区域定位，计算效率提升约50倍。
%
% 2. 设计了自适应遮蔽攻击算法
%    通过逐步增加遮蔽区域数量和半径，并在每个样本上独立进行
%    提前终止判断，实现了以最小扰动面积进行有效攻击的目标。
%    实验表明，自适应遮蔽攻击（N=5,R=3）将标准模型准确率
%    从99.0%降至34.33%，攻击效果与PGD攻击（33.34%）相当。
%
% 3. 提出了混合对抗性训练策略
%    将遮蔽攻击与PGD攻击按比例结合进行对抗性训练，
%    使模型对多种攻击类型均具有防御能力。
%    实验结果显示，混合训练策略在保持98.9%清洁准确率的同时，
%    对PGD攻击的防御准确率达到91.48%，
%    对自适应遮蔽攻击的防御准确率达到76.06%。
%
% 4. 进行了系统的实验评估与分析
%    在MNIST数据集上对10种防御策略和12种攻击类型
%    进行了全面对比实验，揭示了以下关键规律：
%    - 针对单一攻击类型的防御训练对其他攻击类型几乎无效
%    - 自适应遮蔽训练在遮蔽攻击防御中效果最佳
%    - 混合训练是多攻击场景下最实用的防御方案
%    - 参数敏感性分析为遮蔽攻击参数选择提供了实验依据


\section{研究局限性}

% 本文研究存在以下局限性：
%
% 1. 数据集局限
%    实验仅在MNIST数据集上进行验证。
%    MNIST图像分辨率低（28×28）、通道数少（灰度），
%    在更复杂的数据集（如CIFAR-10、ImageNet）上的效果有待验证。
%
% 2. 模型规模局限
%    实验使用的LeNet-5为小规模网络，
%    在大规模模型（如ResNet、VGG等）上的适用性需要进一步研究。
%
% 3. 攻击可见性
%    遮蔽攻击产生的扰动在视觉上较为明显，
%    不同于Lp范数攻击的不可察觉特性，
%    在某些需要隐蔽攻击的场景下适用性有限。
%
% 4. 混合训练的计算开销
%    混合训练需要同时生成两种类型的对抗样本，
%    训练时间约为标准训练的3-4倍。


\section{未来研究方向}

% 基于本文研究，未来可以从以下方向进行深入探索：
%
% 1. 扩展到更复杂的数据集和模型
%    在CIFAR-10、CIFAR-100、ImageNet等数据集上验证方法的有效性，
%    并在ResNet、VGG、WideResNet等深层模型上进行测试。
%
% 2. 改进自适应遮蔽攻击算法
%    探索更高效的区域搜索策略（如二分搜索），
%    减少自适应攻击的计算开销，
%    研究考虑语义信息的遮蔽位置选择。
%
% 3. 多攻击类型的联合防御
%    研究更多攻击类型（如旋转、模糊、颜色变换等）的联合训练，
%    探索最优的攻击类型组合和训练比例。
%
% 4. 理论分析
%    从理论上分析遮蔽攻击与Lp范数攻击之间的关系，
%    研究混合训练策略的收敛性保证和最优比例选择。
%
% 5. 实际应用
%    将遮蔽攻击防御应用于自动驾驶、人脸识别等安全关键场景，
%    研究物理世界中遮蔽攻击的防御方案。
